\documentclass[12pt,a4paper]{article}

% ============================================================
% PISCES Trial - Fish Oil and Cardiovascular Events in Hemodialysis
% 作者: 謝慕揚醫師 MD, PhD, FESC
% ============================================================

% Chinese font setup
\usepackage{xeCJK}
\usepackage{fontspec}
\IfFontExistsTF{Noto Serif CJK TC}{
  \setCJKmainfont{Noto Serif CJK TC}
  \setCJKsansfont{Noto Serif CJK TC}
  \setCJKmonofont{Noto Serif CJK TC}
}{\setCJKmainfont{SimSun}
  \setCJKsansfont{SimSun}
  \setCJKmonofont{SimSun}
}

% Essential packages
\usepackage[margin=2cm]{geometry}
\usepackage{graphicx}
\usepackage{booktabs}
\usepackage{longtable}
\usepackage{array}
\usepackage{xcolor}
\usepackage{hyperref}
\usepackage{tcolorbox}
\usepackage{titlesec}
\usepackage{enumitem}
\usepackage{amsmath}
\usepackage{multirow}
\usepackage{fancyhdr}

% Color definitions
\definecolor{emergency_red}{RGB}{186,24,27}
\definecolor{emergency_blue}{RGB}{0,114,188}
\definecolor{emergency_gray}{RGB}{120,120,120}
\definecolor{highlight_yellow}{RGB}{255,250,205}

% Hyperlink setup
\hypersetup{
    colorlinks=true,
    linkcolor=emergency_blue,
    citecolor=emergency_blue,
    urlcolor=emergency_blue,
    pdfborder={0 0 0}
}

% Title formats
\titleformat{\section}{\Large\bfseries\color{emergency_red}}{\thesection}{1em}{}
\titleformat{\subsection}{\large\bfseries\color{emergency_blue}}{\thesubsection}{1em}{}
\titleformat{\subsubsection}{\normalsize\bfseries\color{emergency_gray}}{\thesubsubsection}{1em}{}

% Header/Footer setup
\pagestyle{fancy}
\fancyhf{}
\fancyhead[L]{PISCES Trial: 魚油與血液透析患者心血管事件}
\fancyhead[R]{\thepage}
\renewcommand{\headrulewidth}{0.4pt}

% ============================================================
% 文件資訊
% ============================================================

\title{\textbf{Fish-Oil Supplementation and Cardiovascular Events in Patients Receiving Hemodialysis\\
魚油補充劑與血液透析患者心血管事件：PISCES 試驗}}
\author{謝慕揚醫師 MD, PhD, FESC\\
資料來源：\href{https://doi.org/10.1056/NEJMoa2513032}{New England Journal of Medicine (2025)}}
\date{整理日期：2025年11月}

\begin{document}

\maketitle

% Key findings box
\begin{tcolorbox}[colback=yellow!10!white,colframe=orange!75!black,title=核心發現摘要]
\textbf{重大發現：}在接受維持性血液透析的患者中，每日補充 n-3 polyunsaturated fatty acids（魚油 4g/天，含 EPA 1.6g + DHA 0.8g）可顯著降低嚴重心血管事件發生率達 43\%（HR 0.57, 95\% CI 0.47-0.70, P<0.001）。

\textbf{臨床意義：}這是首個大型隨機對照試驗證實魚油補充劑對血液透析患者心血管事件的預防效果。考量透析患者心血管死亡率是一般人群的 20 倍，且目前缺乏有效的預防療法，本研究結果具有重大臨床價值。
\end{tcolorbox}

\tableofcontents
\newpage

% ============================================================
\section{研究概述}
% ============================================================

PISCES（Protection against Incidences of Serious Cardiovascular Events Study）是一項多中心、雙盲、隨機、安慰劑對照試驗，評估每日魚油補充劑對接受血液透析患者心血管事件的影響。研究於加拿大與澳洲 26 個中心進行，共納入 1228 位受試者，追蹤 3.5 年。結果顯示魚油組嚴重心血管事件顯著減少，且各項次要終點均呈一致的保護效果。

% ============================================================
\section{研究資訊}
% ============================================================

\subsection{基本資料}

\textbf{標題：}Fish-Oil Supplementation and Cardiovascular Events in Patients Receiving Hemodialysis\\
\textbf{作者：}Lok CE, Farkouh M, Hemmelgarn BR, et al. (for the PISCES Investigators)\\
\textbf{機構：}University Health Network, Toronto, Canada（主導中心）\\
\textbf{期刊：}New England Journal of Medicine (2025)\\
\textbf{DOI：}\href{https://doi.org/10.1056/NEJMoa2513032}{10.1056/NEJMoa2513032}\\
\textbf{研究註冊：}ISRCTN00691795

\subsection{研究背景}

\begin{itemize}
    \item 全球超過 380 萬人接受腎臟替代療法治療末期腎病，其中大多數為血液透析
    \item 心血管疾病影響超過三分之二的血液透析患者，並造成超過 75\% 的相關死亡
    \item 血液透析患者的心血管死亡率是一般人群的 \textbf{20 倍}
    \item 這些高風險患者同時面臨傳統與非傳統心血管危險因子
    \item 目前缺乏經證實有效的心血管事件預防醫療介入措施
    \item n-3 polyunsaturated fatty acids（n-3 PUFA）在一般人群可能有心血管保護效果
    \item 血液透析患者血中 n-3 fatty acid 濃度較一般人群低
\end{itemize}

\subsection{研究假說}

在接受維持性血液透析的患者中，口服補充 EPA 與 DHA（魚油中的長鏈 n-3 PUFA）可降低心血管事件發生率。

% ============================================================
\section{研究方法}
% ============================================================

\subsection{研究設計}

\begin{itemize}
    \item \textbf{研究類型：}多中心、平行組、雙盲、隨機、安慰劑對照試驗
    \item \textbf{研究地點：}加拿大與澳洲共 16 個主要中心、10 個衛星中心
    \item \textbf{收案期間：}2013年11月28日至2019年7月22日
    \item \textbf{追蹤期間：}最長 3.5 年（最後追蹤日期 2023年3月31日）
    \item \textbf{隨機分配：}1:1 比例分配至魚油組或安慰劑組
    \item \textbf{分層因子：}試驗中心、是否有心血管事件病史
\end{itemize}

\subsection{納入與排除條件}

\begin{longtable}{p{3cm}p{11cm}}
\toprule
\textbf{類型} & \textbf{標準} \\
\midrule
\endfirsthead

\toprule
\textbf{類型} & \textbf{標準} \\
\midrule
\endhead

\bottomrule
\endfoot

\textbf{納入條件} &
• 年齡 $\geq$ 18 歲\newline
• 末期腎病（end-stage kidney disease）\newline
• 每週接受血液透析 3-4 次\newline
• 納入前臨床狀況穩定 \\
\midrule
\textbf{排除條件} &
• 隨機分配時正在服用 n-3 fatty acid 補充劑\newline
• 對魚、大豆、玉米或其相關產品過敏 \\
\bottomrule
\end{longtable}

\subsection{介入措施}

\begin{longtable}{p{3cm}p{11cm}}
\toprule
\textbf{組別} & \textbf{介入內容} \\
\midrule
\endfirsthead

\bottomrule
\endfoot

\textbf{魚油組} &
每日口服 4g n-3 PUFA（四顆 1g 膠囊）\newline
• 蒸氣除臭、柑橘口味\newline
• 40:20 ethyl ester 配方\newline
• 每日總劑量：\textbf{EPA 1.6g + DHA 0.8g} \\
\midrule
\textbf{安慰劑組} &
每日口服柑橘口味玉米油膠囊 \\
\bottomrule
\end{longtable}

\subsection{研究終點}

\subsubsection{主要終點（Primary Endpoint）}

所有嚴重心血管事件的複合終點，包括：
\begin{itemize}
    \item 心血管死亡（sudden cardiac death、non-sudden cardiac death）
    \item 致命與非致命 myocardial infarction
    \item 導致截肢的 peripheral vascular disease
    \item 致命與非致命 stroke
\end{itemize}

\textbf{注意：}心臟衰竭未納入主要終點，因血液透析患者常有非心因性的液體過載。

\subsubsection{次要終點（Secondary Endpoints）}

\begin{itemize}
    \item 主要終點擴展至包含非心因性死亡
    \item 主要終點各組成項目
    \item 首次心血管事件或任何原因死亡
\end{itemize}

\subsection{統計分析}

\begin{itemize}
    \item \textbf{分析原則：}Intention-to-treat analysis
    \item \textbf{復發事件分析：}Prentice-Williams-Peterson gap-time model（Cox 比例風險模型的延伸）
    \item \textbf{單一事件分析：}Cox proportional-hazards model
    \item \textbf{檢定力計算：}1100 人追蹤 3.5 年可提供 82\% 檢定力，在 $\alpha$ = 0.05 下偵測 HR 0.825
    \item \textbf{統計軟體：}R software version 4.4.2
\end{itemize}

% ============================================================
\section{主要結果}
% ============================================================

\subsection{受試者特徵}

共 1228 位受試者完成隨機分配並至少接受一劑試驗藥物：
\begin{itemize}
    \item 魚油組：610 人（1394 patient-years 追蹤）
    \item 安慰劑組：618 人（1382 patient-years 追蹤）
\end{itemize}

\begin{longtable}{p{6cm}cc}
\caption{基線患者特徵} \\
\toprule
\textbf{特徵} & \textbf{魚油組 (n=610)} & \textbf{安慰劑組 (n=618)} \\
\midrule
\endfirsthead

\multicolumn{3}{c}{\textbf{表格續接}} \\
\toprule
\textbf{特徵} & \textbf{魚油組} & \textbf{安慰劑組} \\
\midrule
\endhead

\bottomrule
\endfoot

\multicolumn{3}{l}{\textbf{人口學資料}} \\
年齡 (歲) & 64.1±13.5 & 64.5±13.8 \\
男性 & 377 (62.2\%) & 387 (63.1\%) \\
BMI (kg/m²) & 27.3±6.4 & 27.5±6.6 \\
\midrule
\multicolumn{3}{l}{\textbf{種族}} \\
白人 & 244 (40.0\%) & 244 (39.5\%) \\
亞裔 & 99 (16.2\%) & 95 (15.4\%) \\
東南亞裔 & 76 (12.5\%) & 99 (16.0\%) \\
黑人 & 86 (14.1\%) & 71 (11.5\%) \\
\midrule
\multicolumn{3}{l}{\textbf{共病症}} \\
糖尿病 & 342 (56.1\%) & 330 (53.4\%) \\
高血壓 & 522 (85.6\%) & 516 (83.5\%) \\
無心血管疾病史 & 392 (64.3\%) & 402 (65.0\%) \\
\midrule
透析時間中位數 (IQR)，年 & 2.3 (1.1-4.9) & 2.6 (1.1-4.8) \\
\midrule
\multicolumn{3}{l}{\textbf{實驗室數值中位數 (IQR)}} \\
Hemoglobin (g/L) & 109 (101-118) & 110 (102-118) \\
Total cholesterol (mmol/L) & 3.41 (2.79-4.20) & 3.29 (2.84-4.12) \\
LDL cholesterol (mmol/L) & 1.59 (1.15-2.24) & 1.50 (1.10-2.22) \\
HDL cholesterol (mmol/L) & 1.04 (0.83-1.29) & 1.05 (0.84-1.33) \\
Triglycerides (mmol/L) & 1.29 (0.94-1.95) & 1.35 (0.92-1.97) \\
\midrule
\multicolumn{3}{l}{\textbf{藥物使用}} \\
Statin & 332 (54.4\%) & 357 (57.8\%) \\
Beta-blocker & 301 (49.3\%) & 303 (49.0\%) \\
Calcium-channel blocker & 277 (45.4\%) & 251 (40.6\%) \\
RAAS inhibitor & 236 (38.7\%) & 244 (39.5\%) \\
Diuretic & 172 (28.2\%) & 169 (27.3\%) \\
Anticoagulant/antiplatelet & 141 (23.1\%) & 143 (23.1\%) \\
\bottomrule
\end{longtable}

\subsection{主要終點結果}

\begin{tcolorbox}[
    colback=emergency_red!5!white,
    colframe=emergency_red,
    title=\textbf{主要終點結果},
    fonttitle=\bfseries
]
\textbf{嚴重心血管事件發生率：}
\begin{itemize}
    \item 魚油組：0.31 per 1000 patient-days（158 事件）
    \item 安慰劑組：0.61 per 1000 patient-days（309 事件）
    \item \textbf{Hazard Ratio：0.57（95\% CI 0.47-0.70, P<0.001）}
    \item 相對風險降低：\textbf{43\%}
\end{itemize}
\end{tcolorbox}

發生至少一次心血管事件的比例：
\begin{itemize}
    \item 魚油組：20.8\%（127/610）
    \item 安慰劑組：33.7\%（208/618）
\end{itemize}

復發事件（$\geq$2 次）：
\begin{itemize}
    \item 魚油組：4.3\%（26/610）
    \item 安慰劑組：10.5\%（65/618）
\end{itemize}

\subsection{次要終點結果}

\begin{longtable}{p{5.5cm}cccc}
\caption{主要與次要終點分析} \\
\toprule
\textbf{終點} & \textbf{魚油組} & \textbf{安慰劑組} & \textbf{HR} & \textbf{95\% CI} \\
 & \textbf{(事件數)} & \textbf{(事件數)} & & \\
\midrule
\endfirsthead

\toprule
\textbf{終點} & \textbf{魚油組} & \textbf{安慰劑組} & \textbf{HR} & \textbf{95\% CI} \\
\midrule
\endhead

\bottomrule
\endfoot

\multicolumn{5}{l}{\textbf{主要終點}} \\
所有嚴重心血管事件 & 158 & 309 & \textbf{0.57} & 0.47-0.70 \\
\midrule
\multicolumn{5}{l}{\textbf{亞組分析（依心血管病史）}} \\
有心血管事件病史 & 81 & 164 & 0.50 & 0.37-0.67 \\
無心血管事件病史 & 77 & 145 & 0.55 & 0.40-0.76 \\
\midrule
\multicolumn{5}{l}{\textbf{次要終點}} \\
主要終點 + 非心因死亡 & 266 & 381 & 0.77 & 0.65-0.90 \\
任何原因死亡 & 175 & 195 & 0.89 & 0.73-1.01 \\
\midrule
\multicolumn{5}{l}{\textbf{主要終點組成項目}} \\
Cardiac death & 63 & 113 & \textbf{0.55} & 0.40-0.75 \\
Myocardial infarction（致命/非致命） & 49 & 96 & \textbf{0.56} & 0.40-0.80 \\
Peripheral vascular disease $\rightarrow$ amputation & 35 & 66 & \textbf{0.57} & 0.38-0.86 \\
Stroke（致命/非致命） & 11 & 34 & \textbf{0.37} & 0.18-0.76 \\
\midrule
首次心血管事件或任何原因死亡 & 215 & 270 & 0.73 & 0.61-0.87 \\
\bottomrule
\end{longtable}

\subsection{各終點發生比例}

\begin{longtable}{p{6cm}cc}
\toprule
\textbf{心血管事件類型} & \textbf{魚油組} & \textbf{安慰劑組} \\
\midrule
\endfirsthead

\bottomrule
\endfoot

Cardiac death & 10.3\% & 18.3\% \\
Myocardial infarction（致命/非致命） & 6.9\% & 11.7\% \\
Peripheral vascular disease $\rightarrow$ amputation & 4.9\% & 7.1\% \\
Stroke（致命/非致命） & 1.6\% & 4.5\% \\
\bottomrule
\end{longtable}

\subsection{藥物遵從性驗證}

於 3 個月時測量血漿 phospholipid 中 n-3 fatty acid 濃度：
\begin{itemize}
    \item 魚油組 EPA 增加：1.2±1.6（較基線）
    \item 安慰劑組 EPA 變化：0.0±0.5（較基線）
    \item 組間差異：1.3（95\% CI 1.0-1.6）
\end{itemize}
此結果確認受試者遵從用藥，且魚油足以改變循環脂肪酸組成。

\subsection{安全性分析}

\begin{longtable}{p{6cm}cc}
\caption{安全性與嚴重不良事件} \\
\toprule
\textbf{嚴重不良事件} & \textbf{魚油組 (n=610)} & \textbf{安慰劑組 (n=618)} \\
\midrule
\endfirsthead

\bottomrule
\endfoot

感染 & 124 (20.3\%) & 110 (17.8\%) \\
血管通路相關事件 & 51 (8.4\%) & 63 (10.2\%) \\
透析相關事件（液體過載） & 52 (8.5\%) & 48 (7.8\%) \\
腸胃道相關事件 & 56 (9.2\%) & 61 (9.9\%) \\
骨科相關事件 & 36 (5.9\%) & 34 (5.5\%) \\
呼吸系統相關事件 & 31 (5.1\%) & 23 (3.7\%) \\
神經學相關事件 & 23 (3.8\%) & 21 (3.4\%) \\
\midrule
\multicolumn{3}{l}{\textbf{出血事件}} \\
腸胃道出血 & 16 (2.6\%) & 26 (4.2\%) \\
腦出血 & 10 (1.6\%) & 9 (1.5\%) \\
其他出血 & 6 (1.0\%) & 13 (2.1\%) \\
\textbf{出血總計} & \textbf{29 (4.8\%)} & \textbf{47 (7.6\%)} \\
\midrule
無嚴重不良事件 & 310 (50.8\%) & 316 (51.1\%) \\
\bottomrule
\end{longtable}

\begin{tcolorbox}[
    colback=emergency_blue!5!white,
    colframe=emergency_blue,
    title=\textbf{安全性重點},
    fonttitle=\bfseries
]
\begin{itemize}
    \item 兩組嚴重不良事件發生率相似
    \item \textbf{魚油組出血事件反而較少}（4.8\% vs. 7.6\%）
    \item 這與過去認為 n-3 fatty acid 可能增加出血風險的擔憂不符
    \item 多項研究已顯示 n-3 fatty acid 反而降低出血風險
\end{itemize}
\end{tcolorbox}

% ============================================================
\section{討論與臨床意義}
% ============================================================

\subsection{主要發現}

\begin{enumerate}
    \item \textcolor{blue}{\textbf{顯著心血管保護效果：}}每日補充 4g n-3 PUFA 可降低約 40\% 嚴重心血管事件（HR 0.57, P<0.001）

    \item \textcolor{blue}{\textbf{一致性保護效果：}}所有次要終點方向一致
    \begin{itemize}
        \item Cardiac death：HR 0.55
        \item Myocardial infarction：HR 0.56
        \item Peripheral vascular disease $\rightarrow$ amputation：HR 0.57
        \item Stroke：HR 0.37（降低 63\%）
    \end{itemize}

    \item \textcolor{blue}{\textbf{一級與二級預防均有效：}}
    \begin{itemize}
        \item 有心血管病史者：HR 0.50
        \item 無心血管病史者：HR 0.55
    \end{itemize}

    \item \textcolor{blue}{\textbf{安全性良好：}}魚油組出血事件反而較少
\end{enumerate}

\subsection{與過去研究的比較}

\begin{itemize}
    \item 本研究結果與先前 meta-analysis（7 試驗、1045 人）一致
    \item 與 FISH trial 的次要終點結果一致（HR 0.43, P=0.04）
    \item 近期觀察性研究顯示血中高 n-3 fatty acid 濃度與較低心血管風險相關（HR 0.60）
\end{itemize}

\subsection{可能的保護機轉}

n-3 fatty acids（EPA 與 DHA）的潛在心血管保護機轉：

\begin{longtable}{p{4cm}p{10cm}}
\toprule
\textbf{機轉類型} & \textbf{說明} \\
\midrule
\endfirsthead

\bottomrule
\endfoot

抗血栓作用 & 抑制血小板聚集，降低血栓形成風險 \\
抗發炎作用 & 減少 proinflammatory cytokines，改善透析相關發炎狀態 \\
降血脂作用 & 降低 triglycerides，改善脂質代謝 \\
\textbf{抗心律不整作用} &
• 直接抑制 cardiomyocyte sarcolemma voltage-dependent sodium current\newline
• 抑制 L-type calcium current\newline
• 延長 refractory period\newline
• 抑制 depolarization\newline
• 增加細胞電穩定性 \\
心血管重塑作用 & 改善血管功能與心臟重塑 \\
\bottomrule
\end{longtable}

\subsection{為何透析患者可能特別受益？}

\begin{enumerate}
    \item \textbf{基礎 n-3 fatty acid 濃度低：}透析患者血中 EPA、DHA 濃度明顯低於一般人群
    \item \textbf{特殊代謝環境：}促發炎、促心律不整的環境
    \item \textbf{血液透析本身的影響：}
    \begin{itemize}
        \item 快速液體、電解質變化（鈉、鉀、鈣）
        \item 心臟 stunning 與缺血
        \item 血管功能障礙
    \end{itemize}
    \item \textbf{高心血管風險：}死亡率是一般人群的 20 倍，效果更容易顯現
\end{enumerate}

\subsection{為何過去一般人群研究結果不一致？}

\begin{longtable}{p{4cm}p{10cm}}
\toprule
\textbf{可能因素} & \textbf{說明} \\
\midrule
\endfirsthead

\bottomrule
\endfoot

EPA/DHA 比例不同 & 不同配方的 EPA:DHA 比例差異可能影響效果 \\
劑量差異 & 每日 >1g 似乎更有效 \\
配方類型 &
• REDUCE-IT（正向）：ethyl ester 配方，緩慢釋放\newline
• STRENGTH（負向）：carboxylic acid 配方，快速吸收導致更多腸胃道副作用 \\
基線 n-3 濃度 & 基線濃度已經足夠者可能受益較少 \\
基線心血管風險 & 高風險族群更容易看到效果 \\
\bottomrule
\end{longtable}

\subsection{臨床實務建議}

\begin{tcolorbox}[
    colback=emergency_blue!5!white,
    colframe=emergency_blue,
    title=\textbf{臨床建議},
    fonttitle=\bfseries
]
根據 PISCES 試驗結果，對於接受維持性血液透析的患者：

\begin{enumerate}
    \item \textbf{考慮每日補充魚油：}
    \begin{itemize}
        \item 劑量：4g n-3 PUFA/天
        \item 成分：EPA 1.6g + DHA 0.8g（40:20 配方）
        \item 配方：ethyl ester，蒸氣除臭、柑橘口味可改善耐受性
    \end{itemize}

    \item \textbf{適用對象：}無論有無心血管疾病史均可能受益

    \item \textbf{禁忌症：}
    \begin{itemize}
        \item 對魚、大豆、玉米或其產品過敏
    \end{itemize}

    \item \textbf{安全性監測：}
    \begin{itemize}
        \item 魚油不會增加出血風險
        \item 不需特別監測出血併發症
    \end{itemize}
\end{enumerate}
\end{tcolorbox}

% ============================================================
\section{研究限制}
% ============================================================

\begin{enumerate}
    \item \textbf{僅納入血液透析患者：}無法推論至腹膜透析或腎移植患者

    \item \textbf{轉換治療模式者被 censored：}轉為腹膜透析或接受腎移植者的數據被截斷

    \item \textbf{出血事件未經裁決：}僅依照試驗方案定義，未經獨立裁決

    \item \textbf{僅部分受試者進行 n-3 fatty acid 血液檢測：}因預算限制，僅隨機抽樣 232 人

    \item \textbf{心血管藥物使用比例偏低：}
    \begin{itemize}
        \item Statin 使用率 <60\%
        \item 其他降脂藥物 <10\%
        \item 可能反映透析患者常被排除於心血管藥物試驗之外
    \end{itemize}

    \item \textbf{無法確定是否有受試者自行補充魚油：}但血液檢測顯示安慰劑組 EPA/DHA 濃度無變化

    \item \textbf{結果可能無法推論至非透析患者}
\end{enumerate}

% ============================================================
\section{結論}
% ============================================================

\begin{enumerate}
    \item \textcolor{red}{\textbf{主要結論：}}在接受維持性血液透析的患者中，每日補充 n-3 polyunsaturated fatty acids（4g 魚油膠囊，含 EPA 1.6g + DHA 0.8g）可顯著降低嚴重心血管事件發生率（HR 0.57, 95\% CI 0.47-0.70, P<0.001）

    \item \textcolor{red}{\textbf{一致性效果：}}對各項心血管終點（cardiac death、myocardial infarction、stroke、peripheral vascular disease）均呈現保護效果

    \item \textcolor{red}{\textbf{臨床意義：}}考量血液透析患者極高的心血管風險，且目前缺乏有效預防療法，魚油補充劑可能成為這群患者重要的心血管保護策略
\end{enumerate}

\vspace{1cm}
\textbf{關鍵詞：}n-3 polyunsaturated fatty acids、fish oil、EPA、DHA、hemodialysis、cardiovascular events、end-stage kidney disease

% ============================================================
\section{參考文獻}
% ============================================================

\begin{enumerate}
    \item Lok CE, Farkouh M, Hemmelgarn BR, et al. Fish-Oil Supplementation and Cardiovascular Events in Patients Receiving Hemodialysis. \href{https://doi.org/10.1056/NEJMoa2513032}{\textit{N Engl J Med}. 2025.}

    \item Bhatt DL, Steg PG, Miller M, et al. Cardiovascular Risk Reduction with Icosapent Ethyl for Hypertriglyceridemia. \href{https://doi.org/10.1056/NEJMoa1812792}{\textit{N Engl J Med}. 2019;380(1):11-22.}

    \item Nicholls SJ, Lincoff AM, Garcia M, et al. Effect of High-Dose Omega-3 Fatty Acids vs Corn Oil on Major Adverse Cardiovascular Events in Patients at High Cardiovascular Risk: The STRENGTH Randomized Clinical Trial. \href{https://doi.org/10.1001/jama.2020.22258}{\textit{JAMA}. 2020;324(22):2268-2280.}

    \item Lok CE, Moist L, Hemmelgarn BR, et al. Effect of Fish Oil Supplementation on Graft Patency and Cardiovascular Events Among Patients With New Synthetic Arteriovenous Hemodialysis Grafts: A Randomized Controlled Trial. \href{https://doi.org/10.1001/jama.2012.3059}{\textit{JAMA}. 2012;307(17):1809-1816.}

    \item Saglimbene VM, Wong G, van Zwieten A, et al. Effects of Omega-3 Polyunsaturated Fatty Acid Intake in Patients with Chronic Kidney Disease: Systematic Review and Meta-Analysis of Randomized Controlled Trials. \href{https://doi.org/10.1016/j.clnu.2019.02.041}{\textit{Clin Nutr}. 2020;39(2):358-368.}
\end{enumerate}

% ============================================================
\section{縮寫對照表}
% ============================================================

\begin{longtable}{p{4cm}p{10cm}}
\toprule
\textbf{縮寫} & \textbf{全名} \\
\midrule
\endfirsthead

\toprule
\textbf{縮寫} & \textbf{全名} \\
\midrule
\endhead

\bottomrule
\endfoot

CI & Confidence Interval（信賴區間） \\
DHA & Docosahexaenoic Acid（二十二碳六烯酸） \\
EPA & Eicosapentaenoic Acid（二十碳五烯酸） \\
HDL & High-Density Lipoprotein（高密度脂蛋白） \\
HR & Hazard Ratio（風險比） \\
IQR & Interquartile Range（四分位距） \\
LDL & Low-Density Lipoprotein（低密度脂蛋白） \\
n-3 PUFA & n-3 Polyunsaturated Fatty Acids（n-3 多元不飽和脂肪酸） \\
PISCES & Protection against Incidences of Serious Cardiovascular Events Study \\
RAAS & Renin-Angiotensin-Aldosterone System（腎素-血管張力素-醛固酮系統） \\
\bottomrule
\end{longtable}

\end{document}
